\chapter{等价}
\label{cha:equivalences}

现在我们要更详细地研究\emph{类型的等价（equivalence of types）}这一概念，该概念在 \cref{sec:basics-equivalences} 中已有简要介绍。
具体来说，我们将给出几种不同的方式来定义一个类型 $\isequiv(f)$，使其满足那里提到的性质。
回想一下，我们希望 $\isequiv(f)$ 具有以下性质，现重述如下：

\begin{enumerate}
\item $\qinv(f) \to \isequiv (f)$。\label{item:beb1}
\item $\isequiv (f) \to \qinv(f)$。\label{item:beb2}
\item $\isequiv(f)$ 是一个命题。\label{item:beb3}
\end{enumerate}

这里 $\qinv(f)$ 表示 $f$ 的拟逆所构成的类型：
\begin{equation*}
  \sm{g:B\to A} \big((f \circ g \htpy \idfunc[B]) \times (g\circ f \htpy \idfunc[A])\big).
\end{equation*}
利用函数外延性可知，$\qinv(f)$ 等价于类型
\begin{equation*}
  \sm{g:B\to A} \big((f \circ g = \idfunc[B]) \times (g\circ f = \idfunc[A])\big).
\end{equation*}

我们将定义三种不同的类型，它们都满足性质~\ref{item:beb1}--\ref{item:beb3}，分别称为：

\begin{itemize}
\item 半伴随等价，
\item 双可逆映射，
  \index{function!bi-invertible}
  以及
\item 可缩函数。
\end{itemize}

我们还将证明所有这些类型都是彼此等价的。
这些名称故意取得有些冗长，因为一旦我们知道它们彼此等价并且满足性质~\ref{item:beb1}--\ref{item:beb3}，
我们便会直接使用"等价"一词，而无需指明具体选择了哪一种定义。
不过，为了本章进行比较的需要，我们仍然需要为每种定义保留不同的名称。

然而，在具体考察不同的等价概念之前，我们先略微解释一下，为什么需要一个不同于拟逆性的概念。

\section{拟逆}
\label{sec:quasi-inverses}

\index{quasi-inverse|(}%
我们曾说过，$\qinv(f)$ 并不令人满意，因为它不是一个命题；而我们更希望一个给定的函数最多只能以一种方式``成为等价''。
然而，我们尚未给出任何证据表明 $\qinv(f)$ 不是一个命题。
在本节中，我们将展示一个具体的反例。

\begin{lem}\label{lem:qinv-autohtpy}
  如果 $f:A\to B$ 满足 $\qinv (f)$ 是有实例的，那么
  \[\eqv{\qinv(f)}{\Parens{\prd{x:A}(x=x)}}.\]
\end{lem}

\begin{proof}
  根据假设，$f$ 是一个等价；也就是说，我们有 $e:\isequiv(f)$，从而 $(f,e):\eqv A B$。
  由泛等性，$\idtoeqv:(A=B) \to (\eqv A B)$ 是一个等价，因此我们可以假设 $(f,e)$ 具有 $\idtoeqv(p)$ 的形式，其中 $p:A=B$。
  然后利用道路归纳，可以假设 $p$ 就是 $\refl{A}$，此时 $f$ 即为 $\idfunc[A]$。
  这样，我们只需证明 $\eqv{\qinv(\idfunc[A])}{(\prd{x:A}(x=x))}$。
  现在根据定义，我们有
  \[ \qinv(\idfunc[A]) \jdeq
  \sm{g:A\to A} \big((g \htpy \idfunc[A]) \times (g \htpy \idfunc[A])\big).
  \]
  由函数外延性，这等价于
  \[ \sm{g:A\to A} \big((g = \idfunc[A]) \times (g = \idfunc[A])\big).
  \]
  并且根据 \cref{ex:sigma-assoc}，上式又等价于
  \[ \sm{h:\sm{g:A\to A} (g = \idfunc[A])} (\proj1(h) = \idfunc[A])
  \]
  然而，根据 \cref{thm:contr-paths}，$\sm{g:A\to A} (g = \idfunc[A])$ 是以 $(\idfunc[A],\refl{\idfunc[A]})$ 为中心的可缩类型；因此由 \cref{thm:omit-contr} 可知，这个类型等价于 $\idfunc[A] = \idfunc[A]$。
  而再次利用函数外延性，$\idfunc[A] = \idfunc[A]$ 等价于 $\prd{x:A} x=x$。
\end{proof}

\noindent
我们指出，\cref{ex:qinv-autohtpy-no-univalence} 要求给出上述引理的一个避免使用泛等性的证明。

因此，我们需要的是一个类型 $A$，使得 $\prd{x:A}(x=x)$ 有非平凡的元素。
将 $A$ 视为一个高阶群胚，那么 $\prd{x:A}(x=x)$ 的一个实例就是从 $A$ 的恒等函子到其自身的一个自然变换\index{natural!transformation}。
这类变换被称为构成\define{范畴的中心（center of a category）}，
\index{center!of a category}%
\index{category!center of}%
因为自然性公理要求它们与所有态射交换。
在经典情形下，如果 $A$ 就是一个被视作单对象群胚的群，那么这给出的恰好就是通常群论意义上的中心。
这为下面的引理提供了一些动机。

\begin{lem}\label{lem:autohtpy}
  假设我们有一个类型 $A$，其中 $a:A$ 且 $q:a=a$，满足以下条件：
  \begin{enumerate}
  \item 类型 $a=a$ 是一个集合。\label{item:autohtpy1}
  \item 对所有 $x:A$，我们有 $\brck{a=x}$。\label{item:autohtpy2}
  \item 对所有 $p:a=a$，我们有 $p\ct q = q \ct p$。\label{item:autohtpy3}
  \end{enumerate}
  那么存在 $f:\prd{x:A} (x=x)$ 使得 $f(a)=q$。
\end{lem}

\begin{proof}
  设 $g:\prd{x:A} \brck{a=x}$ 如~\ref{item:autohtpy2} 所给出。
  首先我们观察到，每个类型 $\id[A]xy$ 都是一个集合。
  因为"是集合"是一个命题，所以我们可以应用命题截断的归纳原理，并假设 $g(x)=\bproj p$ 且 $g(y)=\bproj{p'}$，其中 $p:a=x$ 且 $p':a=y$。
  在这种情况下，与 $p$ 和 $\opp{p'}$ 复合，便得到一个等价 $\eqv{(x=y)}{(a=a)}$。
  而 $(a=a)$ 根据~\ref{item:autohtpy1} 是一个集合，所以 $(x=y)$ 也是一个集合。

  现在，我们本想通过给每个 $x$ 指派路径 $\opp{g(x)}
  \ct q \ct g(x)$ 来定义 $f$，但这行不通。
  因为 $g(x)$ 并不是 $a=x$ 的元素，而是 $\brck{a=x}$ 的元素；并且类型 $(x=x)$ 可能不是一个命题，所以我们不能对命题截断使用归纳原理。
  我们可以转而应用 \cref{sec:unique-choice} 中提到的技巧：通过唯一性来刻画我们希望构造的对象。
  对每个 $x:A$，定义类型
  \[ B(x) \defeq \sm{r:x=x} \prd{s:a=x} (r = \opp s \ct q\ct s).\]
  我们断言，对于每个 $x:A$，$B(x)$ 都是一个命题。
  由于这个断言本身也是一个命题，我们可以再次对截断应用归纳原理，并假设 $g(x) = \bproj p$，其中 $p:a=x$。
  现在假设给定 $B(x)$ 中的 $(r,h)$ 和 $(r',h')$；那么我们有
  \[ h(p) \ct \opp{h'(p)} : r = r'. \]
  剩下要证明的是，当沿着这个等式传输时，$h$ 与 $h'$ 是等同的。根据在相等类型和函数类型中的传输（\cref{sec:compute-paths,sec:compute-pi}），这归结为要证：对于任何 $s:a=x$，
  \[ h(s) = h(p) \ct \opp{h'(p)} \ct h'(s) \]
  成立。
  但这个等式的两边都是 $(x=x)$ 中元素之间的等式，因此由我们上面的观察（$(x=x)$ 是集合）可知其成立。

  这样，每个 $B(x)$ 都是一个命题；接下来我们断言 $\prd{x:A} B(x)$ 成立。
  给定 $x:A$，我们现在可以调用命题截断的归纳原理，假设 $g(x) = \bproj p$，其中 $p:a=x$。
  我们定义 $r \defeq \opp p \ct q \ct p$；为了构造 $B(x)$ 的实例，剩下要证明对于任何 $s:a=x$，都有 $r = \opp s \ct q \ct s$。
  对道路进行变形，这归结为要证明 $q\ct (p\ct \opp s) = (p\ct \opp s) \ct q$。
  而这正好是~\ref{item:autohtpy3} 的一个实例。
\end{proof}

\begin{thm}\label{thm:qinv-notprop}
  存在类型 $A$ 和 $B$ 以及函数 $f:A\to B$，使得 $\qinv(f)$ 不是一个命题。
\end{thm}

\begin{proof}
  我们只需展示一个类型 $X$，使得 $\prd{x:X} (x=x)$ 不是一个命题即可。
  如同在 \cref{thm:no-higher-ac} 的证明中那样定义 $X\defeq \sm{A:\type} \brck{\bool=A}$。
  此时，我们只需展示一个 $f:\prd{x:X} (x=x)$，使其不等于 $\lam{x} \refl{x}$。

  令 $a \defeq (\bool,\bproj{\refl{\bool}}) : X$，并取 $q:a=a$ 为对应于非恒等的等价 $e:\eqv\bool\bool$ 的道路；这个等价 $e$ 由 $e(\bfalse)\defeq\btrue$ 与 $e(\btrue)\defeq\bfalse$ 定义。
  我们想要应用 \cref{lem:autohtpy} 来构造 $f$。
  根据 $X$ 的定义、子集类型中的相等性（\cref{subsec:prop-subsets}）以及泛等性，我们得到 $\eqv{(a=a)}{(\eqv{\bool}{\bool})}$，这是一个集合，因此条件~\ref{item:autohtpy1} 成立。
  类似地，根据 $X$ 的定义以及子集类型中的相等性，我们得到条件~\ref{item:autohtpy2}。
  最后，根据 \cref{ex:eqvboolbool} 可知，每个等价 $\eqv\bool\bool$ 都等于 $\idfunc[\bool]$ 或 $e$，于是通过分四种情形讨论，我们便能证明~\ref{item:autohtpy3}。

  这样，我们就得到了 $f:\prd{x:X} (x=x)$，满足 $f(a) = q$。
  由于 $e$ 不等于 $\idfunc[\bool]$，因此 $q$ 不等于 $\refl{a}$，从而 $f$ 不等于 $\lam{x} \refl{x}$。
  所以，$\prd{x:X} (x=x)$ 不是一个命题。
\end{proof}

更一般地，\cref{lem:autohtpy} 蕴含了，任何艾伦伯格–麦克兰恩空间 $K(G,1)$（其中 $G$ 是一个非平凡交换\index{group!abelian}群）都将给出一个反例；见 \cref{cha:homotopy}。
我们之前用过的类型 $X$，实际上等价于 $K(\mathbb{Z}_2,1)$。
在 \cref{cha:hits} 我们将看到，圆 $\Sn^1 = K(\mathbb{Z},1)$ 是另一个容易描述的例子。

我们现在转向描述更好的等价概念。

\index{quasi-inverse|)}%

%%%%%%%%%%%%%%%%%%%%%%%%%%%%%%%%%%%%%%
\section{半伴随等价}
\label{sec:hae}
%%%%%%%%%%%%%%%%%%%%%%%%%%%%%%%%%%%%%%

\index{equivalence!half adjoint|(defstyle}%
\index{half adjoint equivalence|(defstyle}%
\index{adjoint!equivalence!of types, half|(defstyle}%

在 \cref{sec:quasi-inverses} 中，我们通过丢弃一个可缩类型的做法，得出了 $\qinv(f)$ 等价于 $\prd{x:A} (x=x)$ 的结论。
粗略来说，类型 $\qinv(f)$ 包含三个数据：$g$、$\eta$ 和 $\epsilon$。
当 $f$ 是一个等价时，其中两个数据（$g$ 和 $\eta$）可以被证明共同构成一个可缩类型。
但问题在于，移除这两个数据后，还会剩下一个数据（$\epsilon$）。
为了解决这个问题，我们的想法是添加一个\emph{额外的}数据，让它与 $\epsilon$ 一起形成一个可缩类型。

\begin{defn}\label{defn:ishae}
  一个函数 $f:A\to B$ 是一个\define{半伴随等价（half adjoint equivalence）}，
  如果存在 $g:B\to A$ 以及同伦 $\eta: g \circ f \htpy \idfunc[A]$ 和 $\epsilon:f \circ g \htpy \idfunc[B]$，
  使得存在一个同伦
  \[\tau : \prd{x:A} \map{f}{\eta x} = \epsilon(fx).\]
\end{defn}

由此我们得到一个类型 $\ishae(f)$，它定义为
\begin{equation*}
  \sm{g:B\to A}{\eta: g \circ f \htpy \idfunc[A]}{\epsilon:f \circ g \htpy \idfunc[B]} \prd{x:A} \map{f}{\eta x} = \epsilon(fx).
\end{equation*}
注意在上述定义中，连接 $\eta$ 和 $\epsilon$ 的相干性\index{coherence}条件只涉及 $f$。
我们也可以转而考虑一个涉及 $g$ 的类似相干性条件：
\[\upsilon : \prd{y:B} \map{g}{\epsilon y} = \eta(gy)\]
并由此得到一个类似的定义 $\ishae'(f)$。

幸运的是，这两个条件互相蕴含：

\begin{lem}\label{lem:coh-equiv}
对于函数 $f : A \to B$ 和 $g:B\to A$，以及同伦 $\eta: g \circ f \htpy \idfunc[A]$ 和 $\epsilon:f \circ g \htpy \idfunc[B]$，以下两个条件逻辑等价：
\begin{itemize}
\item $\prd{x:A} \map{f}{\eta x} = \epsilon(fx)$
\item $\prd{y:B} \map{g}{\epsilon y} = \eta(gy)$
\end{itemize}
\end{lem}

\begin{proof}
只需证明一个方向；另一个方向可以通过将 $A$、$f$ 和 $\eta$ 分别替换为 $B$、$g$ 和 $\epsilon$ 得到。
设 $\tau : \prd{x:A}\;\map{f}{\eta x} = \epsilon(fx)$。
固定 $y : B$。
利用 $\epsilon$ 的自然性并应用 $g$，我们得到如下道路的交换图：
\[\uppercurveobject{{ }}\lowercurveobject{{ }}\twocellhead{{ }}
  \xymatrix@C=3pc{gfgfgy \ar@{=}^-{gfg(\epsilon y)}[r] \ar@{=}_{g(\epsilon (fgy))}[d] & gfgy \ar@{=}^{g(\epsilon y)}[d] \\ gfgy \ar@{=}_{g(\epsilon y)}[r] & gy
  }\]
在图的左侧使用 $\tau(gy)$，我们得到
\[\uppercurveobject{{ }}\lowercurveobject{{ }}\twocellhead{{ }}
  \xymatrix@C=3pc{gfgfgy \ar@{=}^-{gfg(\epsilon y)}[r] \ar@{=}_{gf(\eta (gy))}[d] & gfgy \ar@{=}^{g(\epsilon y)}[d] \\ gfgy \ar@{=}_{g(\epsilon y)}[r] & gy
  }\]
利用 $\eta$ 与 $g \circ f$ 的交换性（\cref{cor:hom-fg}），我们有
\[\uppercurveobject{{ }}\lowercurveobject{{ }}\twocellhead{{ }}
  \xymatrix@C=3pc{gfgfgy \ar@{=}^-{gfg(\epsilon y)}[r] \ar@{=}_{\eta (gfgy)}[d] & gfgy \ar@{=}^{g(\epsilon y)}[d] \\ gfgy \ar@{=}_{g(\epsilon y)}[r] & gy
  }\]
然而，根据 $\eta$ 的自然性，我们还有
\[\uppercurveobject{{ }}\lowercurveobject{{ }}\twocellhead{{ }}
  \xymatrix@C=3pc{gfgfgy \ar@{=}^-{gfg(\epsilon y)}[r] \ar@{=}_{\eta (gfgy)}[d] & gfgy \ar@{=}^{\eta(gy)}[d] \\ gfgy \ar@{=}_{g(\epsilon y)}[r] & gy
  }\]
因此，除右侧同伦外全部消去，便得到 $g(\epsilon y) = \eta(g y)$，如所需。
\end{proof}

但重要的是，我们不在 $\ishae (f)$ 的定义中\emph{同时}包含 $\tau$ 和 $\upsilon$（因此得名"\emph{半}伴随等价"（half adjoint equivalence））。
如果同时包含，那么在消去可缩类型之后，我们仍然会剩下一项数据——除非我们再添加一个更高阶的相干条件。
一般来说，如果在奇数个相干条件之后截断，我们期望会得到一个表现良好的类型。

显然，我们有 $\ishae(f) \to\qinv(f)$：只需忘记相干数据。
反方向则是同伦论与范畴论中一个标准论证的变体。

\begin{thm}\label{thm:equiv-iso-adj}
对任意 $f:A\to B$，我们有 $\qinv(f)\to\ishae(f)$。
\end{thm}

\begin{proof}
假设 $(g,\eta,\epsilon)$ 是 $f$ 的一个拟逆。
我们需要提供一个四元组 $(g',\eta',\epsilon',\tau)$ 来见证 $f$ 是半伴随等价。
要定义 $g'$ 和 $\eta'$，我们可以做出显然的选择：令 $g' \defeq g$，$\eta'\defeq \eta$。
然而，在定义 $\epsilon'$ 时，我们需要考虑 $\tau$ 的构造，因此不能想当然地取 $\epsilon'$ 为 $\epsilon$。
取而代之，我们取
\begin{equation*}
\epsilon'(b) \defeq \opp{\epsilon(f(g(b)))}\ct (\ap{f}{\eta(g(b))}\ct \epsilon(b)).
\end{equation*}
现在我们需要找到
\begin{equation*}
\tau(a): \ap{f}{\eta(a)}=\opp{\epsilon(f(g(f(a))))}\ct (\ap{f}{\eta(g(f(a)))}\ct \epsilon(f(a))).
\end{equation*}
首先注意，由 \cref{cor:hom-fg} 我们有
%$\eta(g(f(a)))\ct\eta(a)=\ap{g}{\ap{f}{\eta(a)}}\ct\eta(a)$ and hence it follows that
$\eta(g(f(a)))=\ap{g}{\ap{f}{\eta(a)}}$。
因此，我们可以应用
\cref{lem:htpy-natural} 来计算
\begin{align*}
\ap{f}{\eta(g(f(a)))}\ct \epsilon(f(a))
& = \ap{f}{\ap{g}{\ap{f}{\eta(a)}}}\ct \epsilon(f(a))\\
& = \epsilon(f(g(f(a))))\ct \ap{f}{\eta(a)}
\end{align*}
由此我们便得到了所需的道路 $\tau(a)$。
\end{proof}

结合 \cref{lem:coh-equiv}（或对证明作对称化处理），我们也有 $\qinv(f)\to\ishae'(f)$。

尚需证明 $\ishae(f)$ 是一个命题。
为此，我们需要先知道等价映射的纤维是可缩的。

\begin{defn}\label{defn:homotopy-fiber}
  一个映射 $f:A\to B$ 在点 $y:B$ 处的 \define{纤维}
  \indexdef{fiber}%
  \indexsee{function!fiber of}{fiber}%
  定义为：
  \[ \hfib f y \defeq \sm{x:A} (f(x) = y).\]
\end{defn}

在同伦论中，这被称为 $f$ 的\emph{同伦纤维（homotopy fiber）}。
\cref{sec:computational} 中的道路引理导出了纤维中道路的如下刻画：

\begin{lem}\label{lem:hfib}
  对于任意 $f : A \to B$，$y : B$，以及 $(x,p),(x',p') : \hfib{f}{y}$，我们有
  \[ \big((x,p) = (x',p')\big) \eqvsym \Parens{\sm{\gamma : x = x'} f(\gamma) \ct p' = p} \qedhere\]
\end{lem}

\begin{thm}\label{thm:contr-hae}
  如果 $f:A\to B$ 是一个半伴随等价，那么对于任意 $y:B$，纤维 $\hfib f y$ 都是可缩的。
\end{thm}

\begin{proof}
  令 $(g,\eta,\epsilon,\tau) : \ishae(f)$，并固定 $y : B$。
  我们选取 $(gy, \epsilon y)$ 作为 $\hfib{f}{y}$ 的收缩中心。
  现在任取 $(x,p) : \hfib{f}{y}$；我们想要构造一条从 $(gy, \epsilon y)$ 到 $(x,p)$ 的道路。
  根据 \cref{lem:hfib}，只需给出一条道路 $\gamma : \id{gy}{x}$ 使得 $\ap f\gamma \ct p = \epsilon y$ 即可。
  我们令 $\gamma \defeq \opp{g(p)} \ct \eta x$。
  于是有
  \begin{align*}
    f(\gamma) \ct p & = \opp{fg(p)} \ct f (\eta x) \ct p \\
    & = \opp{fg(p)} \ct \epsilon(fx) \ct p \\
    & = \epsilon y
  \end{align*}
  其中第二个等式由 $\tau x$ 推出，第三个等式是 $\epsilon$ 的自然性。
\end{proof}

现在我们来定义封装可缩数据对的类型。
下面的类型将拟逆 $g$ 与其中一条同伦组合在一起。

\begin{defn}\label{defn:linv-rinv}
  给定一个函数 $f:A\to B$，我们定义类型
    \begin{align*}
      \linv(f) &\defeq \sm{g:B\to A} (g\circ f\htpy \idfunc[A])\\
      \rinv(f) &\defeq \sm{g:B\to A} (f\circ g\htpy \idfunc[B])
    \end{align*}
  分别为 $f$ 的 \define{左逆}
  \indexdef{left!inverse}%
  \indexdef{inverse!left}%
  与 \define{右逆}
  \indexdef{right!inverse}%
  \indexdef{inverse!right}%
  。
  如果 $\linv(f)$ 有实例，我们就称 $f$ 是 \define{左可逆的}
  \indexdef{function!left invertible}%
  \indexdef{function!right invertible}%
  ；类似地，如果 $\rinv(f)$ 有实例，则称 $f$ 是 \define{右可逆的}
  \indexdef{left!invertible function}%
  \indexdef{right!invertible function}%
  。
\end{defn}

\begin{lem}\label{thm:equiv-compose-equiv}
  如果 $f:A\to B$ 有一个拟逆，那么 \begin{align*}
    (f\circ \blank) &: (C\to A) \to (C\to B)\\
    (\blank\circ f) &: (B\to C) \to (A\to C).
  \end{align*} 亦然。
\end{lem}

\begin{proof}
  如果 $g$ 是 $f$ 的一个拟逆，那么 $(g\circ \blank)$ 和 $(\blank\circ g)$ 就分别是 $(f\circ \blank)$ 和 $(\blank\circ f)$ 的拟逆。
\end{proof}

\begin{lem}\label{lem:inv-hprop}
  如果 $f : A \to B$ 有一个拟逆，那么类型 $\rinv(f)$ 和 $\linv(f)$ 都是可缩的。
\end{lem}

\begin{proof}
  由函数外延性，我们有
  \[\eqv{\linv(f)}{\sm{g:B\to A} (g\circ f = \idfunc[A])}.\]
  但这就是 $(\blank\circ f)$ 在 $\idfunc[A]$ 上的纤维，因此由 \cref{thm:equiv-compose-equiv,thm:equiv-iso-adj,thm:contr-hae}，它是可缩的。
  类似地，$\rinv(f)$ 等价于 $(f\circ \blank)$ 在 $\idfunc[B]$ 上的纤维，因而也是可缩的。
\end{proof}

接下来我们定义那些将另一个同伦与额外的相干数据组合在一起的类型。\index{coherence}%

\begin{defn}\label{defn:lcoh-rcoh}
对于 $f : A \to B$，一个左逆 $(g,\eta) : \linv(f)$，以及一个右逆 $(g,\epsilon) : \rinv(f)$，我们记
\begin{align*}
\lcoh{f}{g}{\eta} & \defeq \sm{\epsilon : f\circ g \htpy \idfunc[B]} \prd{y:B} g(\epsilon y) = \eta (gy), \\
\rcoh{f}{g}{\epsilon} & \defeq \sm{\eta : g\circ f \htpy \idfunc[A]} \prd{x:A} f(\eta x) = \epsilon (fx).
\end{align*}
\end{defn}

\begin{lem}\label{lem:coh-hfib}
对于任何 $f,g,\epsilon,\eta$，我们有
\begin{align*}
\lcoh{f}{g}{\eta} & \eqvsym {\prd{y:B} \id[\hfib{g}{gy}]{(fgy,\eta(gy))}{(y,\refl{gy})}}, \\
\rcoh{f}{g}{\epsilon} & \eqvsym {\prd{x:A} \id[\hfib{f}{fx}]{(gfx,\epsilon(fx))}{(x,\refl{fx})}}.
\end{align*}
\end{lem}

\begin{proof}
利用 \cref{lem:hfib}。
\end{proof}

\begin{lem}\label{lem:coh-hprop}
  如果 $f$ 是一个半伴随等价，那么对于任何 $(g,\epsilon) : \rinv(f)$，类型 $\rcoh{f}{g}{\epsilon}$ 都是可缩的。
\end{lem}

\begin{proof}
  由 \cref{lem:coh-hfib} 以及依值函数类型保持可缩空间这一事实，只需证明对每个 $x:A$，类型 $\id[\hfib{f}{fx}]{(gfx,\epsilon(fx))}{(x,\refl{fx})}$ 是可缩的。
  但由 \cref{thm:contr-hae}，$\hfib{f}{fx}$ 是可缩的，而任何可缩空间的道路空间自身也是可缩的。
\end{proof}

\begin{thm}\label{thm:hae-hprop}
  对于任何 $f : A \to B$，类型 $\ishae(f)$ 是一个命题。
\end{thm}

\begin{proof}
  由 \cref{ex:prop-inhabcontr}，可以假设 $f$ 是一个半伴随等价，并证明 $\ishae(f)$ 是可缩的。
  现在根据 $\Sigma$ 类型的结合律（\cref{ex:sigma-assoc}），类型 $\ishae(f)$ 等价于
  \[\sm{u : \rinv(f)} \rcoh{f}{\proj{1}(u)}{\proj{2}(u)}.\]
  但由 \cref{lem:inv-hprop,lem:coh-hprop} 以及 $\Sigma$ 类型保持可缩性这一事实，后一类型也是可缩的。
\end{proof}

这样，我们就证明了 $\ishae(f)$ 具备了类型 $\isequiv(f)$ 所需的全部三个理想性质。
在接下来的两节中，我们将考虑另外几种可能性。

\index{equivalence!half adjoint|)}%
\index{half adjoint equivalence|)}%
\index{adjoint!equivalence!of types, half|)}%

\section{双可逆映射}
\label{sec:biinv}

\index{function!bi-invertible|(defstyle}%
\index{bi-invertible function|(defstyle}%
\index{equivalence!as bi-invertible function|(defstyle}%

使用 \cref{sec:hae} 中引入的语言，我们可以将 \cref{sec:basics-equivalences} 中提出的定义重述如下。

\begin{defn}\label{defn:biinv}
  我们称 $f:A\to B$ 是\define{双可逆的（bi-invertible）}，
  如果它同时有一个左逆和一个右逆：
  \[ \biinv (f) \defeq \linv(f) \times \rinv(f). \]
\end{defn}

在 \cref{sec:basics-equivalences} 中，我们证明了 $\qinv(f)\to\biinv(f)$ 以及 $\biinv(f)\to\qinv(f)$。
剩下需要证明的是以下定理。

\begin{thm}\label{thm:isprop-biinv}
  对于任何 $f:A\to B$，类型 $\biinv(f)$ 是一个命题。
\end{thm}

\begin{proof}
  可以假设 $f$ 是双可逆的，并证明 $\biinv(f)$ 是可缩的。
  但由于 $\biinv(f)\to\qinv(f)$，根据 \cref{lem:inv-hprop}，在这种情况下 $\linv(f)$ 和 $\rinv(f)$ 都是可缩的，而可缩类型的积是可缩的。
\end{proof}

注意，这也符合我们在 \cref{sec:hae} 开头提出的建议：我们将 $g$ 和 $\eta$ 结合成一个可缩的类型，并添加一项额外数据，该数据再与 $\epsilon$ 结合成另一个可缩的类型。
区别在于，我们不是添加一个\emph{高阶}数据（一条 2 维道路）来与 $\epsilon$ 结合，而是添加了一个\emph{低阶}数据（一个独立于左逆的右逆）。

\begin{cor}\label{thm:equiv-biinv-isequiv}
  对于任何 $f:A\to B$，我们有 $\eqv{\biinv(f)}{\ishae(f)}$。
\end{cor}

\begin{proof}
  我们有 $\biinv(f) \to \qinv(f) \to \ishae(f)$ 和 $\ishae(f) \to \qinv(f) \to \biinv(f)$。
  由于 $\ishae(f)$ 和 $\biinv(f)$ 都是命题，该等价由 \cref{lem:equiv-iff-hprop} 得出。
\end{proof}

\index{function!bi-invertible|)}%
\index{bi-invertible function|)}%
\index{equivalence!as bi-invertible function|)}%

\section{具有可缩纤维的映射}
\label{sec:contrf}

\index{function!contractible|(defstyle}%
\index{contractible!function|(defstyle}%
\index{equivalence!as contractible function|(defstyle}%

注意，我们关于 $\ishae(f)$ 和 $\biinv(f)$ 的证明，本质上使用了"等价映射的纤维是可缩的"这一事实。
事实上，这一性质本身就可以作为等价的充分定义。

\begin{defn}[可缩映射] \label{defn:equivalence}
  一个映射 $f:A\to B$ 被称为\define{可缩的（contractible）}，
  如果对所有 $y:B$，纤维 $\hfib f y$ 都是可缩的。
\end{defn}

因此，类型 $\iscontr(f)$ 定义为
\begin{align}
  \iscontr(f) &\defeq \prd{y:B} \iscontr(\hfib f y)\label{eq:iscontrf}
  % \\
  % &\defeq \prd{y:B} \iscontr (\setof{x:A | f(x) = y}).
\end{align}
注意在 \cref{sec:contractibility} 中我们定义了\emph{类型}可缩的含义。
这里我们定义的是\emph{映射}可缩的含义。
我们的术语遵循同伦论的一般惯例：如果一个映射的所有（同伦）纤维都具有某个性质，我们就说这个映射具有该性质。
因此，一个类型 $A$ 可缩当且仅当映射 $A\to\unit$ 可缩。
从 \cref{cha:hlevels} 开始，我们也将可缩的映射和类型称为\emph{$(-2)$-截断的}。

我们已经在 \cref{thm:contr-hae} 中证明了 $\ishae(f) \to \iscontr(f)$。
反过来：

\begin{thm}\label{thm:lequiv-contr-hae}
对任意 $f:A\to B$，有 ${\iscontr(f)} \to {\ishae(f)}$。
\end{thm}

\begin{proof}
令 $P : \iscontr(f)$。我们如下定义逆映射 $g : B \to A$：将每个 $y : B$ 送到 $y$ 处纤维的收缩中心：
\[ g(y) \defeq \proj{1}(\proj{1}(Py)). \]
于是可以如下定义同伦 $\epsilon$：将 $y$ 映为 $g(y)$ 确实属于 $y$ 处纤维的见证：
\[ \epsilon(y) \defeq \proj{2}(\proj{1}(P y)). \]
剩下还需要定义 $\eta$ 和 $\tau$。这当然相当于给出 $\rcoh{f}{g}{\epsilon}$ 的一个元素。由 \cref{lem:coh-hfib}，这等价于对每个 $x:A$，在 $f$ 在 $fx$ 处的纤维中给出一条从 $(gfx,\epsilon(fx))$ 到 $(x,\refl{fx})$ 的道路。但这一点很容易：对于任意 $x : A$，类型 $\hfib{f}{fx}$
由假设是可缩的，因此这样一条道路必然存在。我们可以显式地将其构造为
\[\opp{\big(\proj{2}(P(fx))(gfx,\epsilon(fx))\big)} \ct \big(\proj{2}(P(fx)) (x,\refl{fx})\big). \qedhere \]
\end{proof}

也很容易看出：

\begin{lem}\label{thm:contr-hprop}
  对任意 $f$，类型 $\iscontr(f)$ 是一个命题。
\end{lem}

\begin{proof}
  由 \cref{thm:isprop-iscontr}，每个类型 $\iscontr (\hfib f y)$ 都是命题。
  因此，由 \cref{thm:isprop-forall}，\eqref{eq:iscontrf} 也是命题。
\end{proof}

\begin{thm}\label{thm:equiv-contr-hae}
  对任意 $f:A\to B$，有 $\eqv{\iscontr(f)}{\ishae(f)}$。
\end{thm}

\begin{proof}
  我们已经建立了逻辑等价 ${\iscontr(f)} \Leftrightarrow {\ishae(f)}$，并且两者都是命题（\cref{thm:contr-hprop,thm:hae-hprop}）。
  因此 \cref{lem:equiv-iff-hprop} 适用。
\end{proof}

通常我们通过展示拟逆来证明一个函数是等价，但有时上述定义更为方便。
例如，它意味着在证明一个函数是等价时，我们可以自由地假设其陪域是有实例的。

\begin{cor}\label{thm:equiv-inhabcod}
  如果 $f:A\to B$ 满足 $B\to \isequiv(f)$，那么 $f$ 是一个等价。
\end{cor}

\begin{proof}
  要证明 $f$ 是等价，只需证明对任意 $y:B$，$\hfib f y$ 是可缩的。
  但如果有 $e:B\to \isequiv(f)$，那么给定任意这样的 $y$，我们有 $e(y):\isequiv(f)$，因此 $f$ 是等价，从而 $\hfib f y$ 可缩，如所欲证。
\end{proof}

\index{function!contractible|)}%
\index{contractible!function|)}%
\index{equivalence!as contractible function|)}%

\section{关于等价的定义}
\label{sec:concluding-remarks}

\indexdef{equivalence}
我们已经证明等价的三种定义都满足三条所期望的性质，并且彼此等价：
\[ \iscontr(f) \eqvsym \ishae(f) \eqvsym \biinv(f). \]
（等价还有更多可能的定义，但我们只考虑这三种。
更多定义可参见 \cref{ex:brck-qinv} 和本章的习题。）
因此，我们可以选择其中任意一种作为 $\isequiv (f)$ 的"那个"定义。
为明确起见，我们选择定义
\[ \isequiv(f) \defeq \ishae(f).\]
\index{mathematics!formalized}%
这一选择对形式化有利，因为 $\ishae(f)$ 包含了最直接有用的数据。
另一方面，对于其他目的，$\biinv(f)$ 往往更容易处理，因为它不包含 $2$-维道路，而且它的两个对称部分可以独立处理。
不过，就本书而言，具体选择哪一种差别不大。

在本章的剩余部分，我们将研究等价的其他一些性质与刻画。
\index{equivalence!properties of}%

\section{满射与嵌入}
\label{sec:mono-surj}

\index{set}
当 $A$ 和 $B$ 都是集合且 $f:A\to B$ 是一个等价时，我们也称它为\define{同构（isomorphism）}
\indexdef{isomorphism!of sets}%
或\define{双射（bijection）}。
\indexdef{bijection}%
\indexsee{function!bijective}{bijection}%
（对于不是集合的类型，我们避免使用这些词，因为在同伦论和高阶范畴论中，它们通常表示比同伦等价更严格的"相同"概念。）
在集合论中，一个函数是双射当且仅当它既是单射又是满射。
在类型论中同样如此，只要我们适当地表述这些条件。
为了清晰起见，在处理不是集合的类型时，我们将用\emph{嵌入（embedding）}一词代替单射。

\begin{defn}\label{defn:surj-emb}
  设 $f:A\to B$。
  \begin{enumerate}
  \item 若对每个 $b:B$ 都有 $\brck{\hfib f b}$，则称 $f$ 是\define{满的（surjective）}
    \indexsee{surjective!function}{function, surjective}%
    \indexdef{function!surjective}%
    （或称\define{满射（surjection）}）
    \indexsee{surjection}{function, surjective}%
    。
  \item 若对每对 $x,y:A$，函数 $\apfunc f : (\id[A]xy) \to (\id[B]{f(x)}{f(y)})$ 都是一个等价，则称 $f$ 是一个\define{嵌入（embedding）}
    \indexdef{function!embedding}%
    \indexsee{embedding}{function, embedding}%
    。
  \end{enumerate}
\end{defn}

换言之，$f$ 是满射当且仅当 $f$ 的每个纤维都仅知有实例，或者说，对于所有 $b:B$，都仅知存在一个 $a:A$ 使得 $f(a)=b$。
用传统逻辑记号来说，$f$ 是满射当 $\fall{b:B}\exis{a:A} (f(a)=b)$。
这必须与更强的断言 $\prd{b:B}\sm{a:A} (f(a)=b)$ 区分开来；若后者成立，则称 $f$ 是一个\define{可分裂满射（split surjection）}。
\indexsee{split!surjection}{function, split surjective}%
\indexsee{surjection!split}{function, split surjective}%
\indexsee{surjective!function!split}{function, split surjective}%
\indexdef{function!split surjective}%
（由于后一个类型等价于 $\sm{g:B\to A}\prd{b:B} (f(g(b))=b)$，是可分裂满射就等于是\emph{收缩}，如 \cref{sec:contractibility} 所定义。）
\index{retraction}%
\index{function!retraction}%

\cref{sec:axiom-choice} 中的选择公理确切地说就是：\emph{集合之间}的每个满射都是可分裂的。
然而，在泛等公理成立的前提下，说\emph{所有}满射都可分裂是根本不成立的。
在 \cref{thm:no-higher-ac} 中，我们构造了一个类型族 $Y:X\to \type$，使得 $\prd{x:X} \brck{Y(x)}$ 但 $\neg \prd{x:X} Y(x)$；
对于任何这样的族，其第一投影 $(\sm{x:X} Y(x)) \to X$ 就是一个不可分裂的满射。

如果 $A$ 和 $B$ 都是集合，那么由 \cref{lem:equiv-iff-hprop}，$f$ 是嵌入当且仅当
\begin{equation}
  \prd{x,y:A} (\id[B]{f(x)}{f(y)}) \to (\id[A]xy).\label{eq:injective}
\end{equation}
在这种情况下，我们说 $f$ 是\define{单的（injective）}，
\indexsee{injective function}{function, injective}%
\indexdef{function!injective}%
或称它是一个\define{单射（injection）}。
\indexsee{injection}{function, injective}%
对于不是集合的类型，我们避免使用这些词，因为它们可能被解读为 \eqref{eq:injective}，这对于非集合来说是一个性质不佳的概念。
同样成立的是，集合之间的函数是满射当且仅当它在某种合适意义下是\emph{满态射（epimorphism）}，但这对于更一般的类型也不成立，并且满射性通常是更重要的概念。

\begin{thm}\label{thm:mono-surj-equiv}
  一个函数 $f:A\to B$ 是等价，当且仅当它同时是满射和嵌入。
\end{thm}

\begin{proof}
  如果 $f$ 是一个等价，那么每个 $\hfib f b$ 都是可缩的，从而 $\brck{\hfib f b}$ 也是可缩的，所以 $f$ 是满射。
  并且我们在 \cref{thm:paths-respects-equiv} 中已经证明过，任何等价都是嵌入。

  反之，假设 $f$ 既是满射又是嵌入。
  令 $b:B$；我们来证明 $\sm{x:A}(f(x)=b)$ 是可缩的。
  因为 $f$ 是满射，仅知存在 $a:A$ 使得 $f(a)=b$。
  因此，$f$ 在 $b$ 处的纤维是有实例的；剩下只需证明它是一个命题。
  为此，假设给定 $x,y:A$ 以及 $p:f(x)=b$ 和 $q:f(y)=b$。
  那么由于 $\apfunc f$ 是一个等价，存在 $r:x=y$ 使得 $\apfunc f (r) = p \ct \opp q$。
  然而，利用 $\Sigma$-类型中道路的表征，后一个等式可重新整理为 $\trans{r}{p} = q$。
  于是，连同 $r$ 一起，它就在 $f$ 在 $b$ 处的纤维中给出了 $(x,p) = (y,q)$。
\end{proof}

\begin{cor}
  对于任意 $f:A\to B$，我们有
  \[ \isequiv(f) \eqvsym (\mathsf{isEmbedding}(f) \times \mathsf{isSurjective}(f)).\]
\end{cor}

\begin{proof}
  "是满射"和"是嵌入"都是命题；现在应用 \cref{lem:equiv-iff-hprop} 即可。
\end{proof}

当然，我们不能把这当作"等价"的定义，因为嵌入的定义本身就依赖于等价。
不过，这一刻画仍然有用；参见 \cref{sec:whitehead}。
我们将在 \cref{cha:hlevels} 中对其加以推广。

% \section{Fiberwise equivalences}
\section{等价的闭包性质}
\label{sec:equiv-closures}
\label{sec:fiberwise-equivalences}
\index{equivalence!properties of}%

% We end this chapter by observing some important closure properties of equivalences.
我们已经在 \cref{thm:equiv-eqrel} 中看到，等价在复合运算下是封闭的。
此外，还有：

\begin{thm}[2-out-of-3 性质]\label{thm:two-out-of-three}
  \index{2-out-of-3 property}%
  设 $f:A\to B$ 和 $g:B\to C$。
  如果 $f$、$g$ 和 $g\circ f$ 三者中有任意两个是等价，那么第三个也是等价。
\end{thm}

\begin{proof}
  若 $g\circ f$ 和 $g$ 是等价，则 $\opp{(g\circ f)} \circ g$ 是 $f$ 的一个拟逆。
  一方面，我们有 $\opp{(g\circ f)} \circ g \circ f \htpy \idfunc[A]$；另一方面，则有
  \begin{align*}
    f \circ \opp{(g\circ f)} \circ g
    &\htpy \opp g \circ g \circ f \circ \opp{(g\circ f)} \circ g\\
    &\htpy \opp g \circ g\\
    &\htpy \idfunc[B].
  \end{align*}
  类似地，若 $g\circ f$ 和 $f$ 是等价，则 $f\circ \opp{(g\circ f)}$ 是 $g$ 的一个拟逆。
\end{proof}

这是同伦论中关于等价的一个标准闭包条件。
同样广为人知的是，等价对收缩核也封闭，具体含义如下。

\index{retract!of a function|(defstyle}%

\begin{defn}\label{defn:retract}
  称一个函数 $g:A\to B$ 为函数 $f:X\to Y$ 的\define{收缩核}，
  如果存在一个图表
  \begin{equation*}
  \xymatrix{
    {A} \ar[r]^{s} \ar[d]_{g}
    &
    {X} \ar[r]^{r} \ar[d]_{f}
    &
    {A} \ar[d]^{g}
    \\
    {B} \ar[r]_{s'}
    &
    {Y} \ar[r]_{r'}
    &
    {B}
  }
\end{equation*}
  并满足以下条件：
  \begin{enumerate}
    \item 一个同伦 $R:r\circ s \htpy \idfunc[A]$。
    \item 一个同伦 $R':r'\circ s' \htpy\idfunc[B]$。
    \item 一个同伦 $L:f\circ s\htpy s'\circ g$。
    \item 一个同伦 $K:g\circ r\htpy r'\circ f$。
    \item 对每个 $a:A$，存在一条道路 $H(a)$，见证如下方块的交换性：
      \begin{equation*}
  \xymatrix@C=3pc{
    {g(r(s(a)))} \ar@{=}[r]^-{K(s(a))} \ar@{=}[d]_{\ap g{R(a)}}
    &
    {r'(f(s(a)))} \ar@{=}[d]^{\ap{r'}{L(a)}}
    \\
    {g(a)} \ar@{=}[r]_-{\opp{R'(g(a))}}
    &
    {r'(s'(g(a)))}
  }
\end{equation*}
  \end{enumerate}
\end{defn}

回顾 \cref{sec:contractibility}，其中我们定义了何为一个类型是另一个类型的收缩核。
那是上述定义在 $B$ 和 $Y$ 均为 $\unit$ 时的特例。
反过来，正如可缩性的情况一样，映射的收缩核会诱导其纤维的收缩核。

\begin{lem}\label{lem:func_retract_to_fiber_retract}
  如果函数 $g:A\to B$ 是函数 $f:X\to Y$ 的收缩核，那么对每个 $b:B$，
  $\hfib{g}b$ 是 $\hfib{f}{s'(b)}$ 的收缩核，
  其中 $s':B\to Y$ 如 \cref{defn:retract} 中所述。
\end{lem}

\begin{proof}
  假设 $g:A\to B$ 是 $f:X\to Y$ 的收缩核。那么对任意 $b:B$，我们有两个函数
  \begin{align*}
\varphi_b &:\hfiber{g}b\to\hfib{f}{s'(b)}, &
\varphi_b(a,p) & \defeq \pairr{s(a),L(a)\ct s'(p)},\\
\psi_b &:\hfib{f}{s'(b)}\to\hfib{g}b, &
\psi_b(x,q) &\defeq \pairr{r(x),K(x)\ct r'(q)\ct R'(b)}.
\end{align*}
  然后我们有 $\psi_b(\varphi_b({a,p}))\equiv\pairr{r(s(a)),K(s(a))\ct r'(L(a)\ct s'(p))\ct R'(b)}$。
  我们断言：对所有 $b:B$，$\psi_b$ 是一个以 $\varphi_b$ 为截面的收缩，
  也就是说，对所有 $(a,p):\hfib g b$，有 $\psi_b(\varphi_b({a,p}))= \pairr{a,p}$。
  换言之，我们希望证明
  \begin{equation*}
\prd{b:B}{a:A}{p:g(a)=b} \psi_b(\varphi_b({a,p}))= \pairr{a,p}.
\end{equation*}
  通过调换前两个 $\Pi$ 的位置并应用 \cref{thm:omit-contr} 的一种形式，上式等价于
  \begin{equation*}
\prd{a:A}\psi_{g(a)}(\varphi_{g(a)}({a,\refl{g(a)}}))=\pairr{a,\refl{g(a)}}.
\end{equation*}
  对任意 $a$，由 \cref{thm:path-sigma}，此配对的相等等价于一对相等。
  第一个分量由 $R(a):r(s(a))= a$ 可知相等，因此我们只需证明
  \begin{equation*}
\trans{R(a)}{K(s(a))\ct r'(L(a))\ct R'(g(a))} = \refl{g(a)}.
\end{equation*}
  但这个输运计算为 $\opp{g(R(a))}\ct K(s(a))\ct r'(L(a))\ct R'(g(a))$，所以所需的道路由 $H(a)$ 给出。
\end{proof}

\begin{thm}\label{thm:retract-equiv}
  如果 $g$ 是某个等价 $f$ 的收缩核，那么 $g$ 也是一个等价。
\end{thm}

\begin{proof}
  由 \cref{lem:func_retract_to_fiber_retract}，$g$ 的每个纤维都是 $f$ 的某个纤维的收缩核。
  因此，由 \cref{thm:retract-contr}，如果 $f$ 的所有纤维都是可缩的，那么 $g$ 的所有纤维也是可缩的。
\end{proof}

\index{retract!of a function|)}%

\index{fibration}%
\index{total!space}%
最后，我们说明逐纤维等价可以用全空间的等价来刻画。
为解释相关术语，回顾 \cref{sec:fibrations}：一个类型族 $P:A\to\type$ 可以看作以 $A$ 为底空间的纤维化，
其全空间为 $\sm{x:A} P(x)$，而该纤维化即投影映射 $\proj1:\sm{x:A} P(x) \to A$。
从这个角度看，给定两个类型族 $P,Q:A\to\type$，我们可以称一个函数 $f:\prd{x:A} (P(x)\to Q(x))$ 为\define{逐纤维映射}或\define{逐纤维变换}。
\indexsee{transformation!fiberwise}{fiberwise transformation}%
\indexsee{function!fiberwise}{fiberwise transformation}%
\index{fiberwise!transformation|(defstyle}%
\indexsee{fiberwise!map}{fiberwise transformation}%
\indexsee{map!fiberwise}{fiberwise transformation}
这样的映射会诱导一个全空间之间的函数：

\begin{defn}\label{defn:total-map}
  给定类型族 $P,Q:A\to\type$ 和一个映射 $f:\prd{x:A} P(x)\to Q(x)$，定义
  \begin{equation*}
    \total f  \defeq \lam{w}\pairr{\proj{1}w,f(\proj{1}w,\proj{2}w)} : \sm{x:A}P(x)\to\sm{x:A}Q(x).
  \end{equation*}
\end{defn}

\begin{thm}\label{fibwise-fiber-total-fiber-equiv}
  设 $f$ 是类型 $A$ 上两个族 $P$ 和 $Q$ 之间的一个逐纤维变换，并令 $x:A$ 和 $v:Q(x)$。
  那么我们有如下等价
  \begin{equation*}
\eqv{\hfib{\total{f}}{\pairr{x,v}}}{\hfib{f(x)}{v}}.
\end{equation*}
\end{thm}

\begin{proof}
  我们计算如下：
  \begin{align}
  \hfib{\total{f}}{\pairr{x,v}}
  & \jdeq \sm{w:\sm{x:A}P(x)}\pairr{\proj{1}w,f(\proj{1}w,\proj{2}w)}=\pairr{x,v}
  \notag \\
  & \eqv{}{} \sm{a:A}{u:P(a)}\pairr{a,f(a,u)}=\pairr{x,v}
  \tag{by~\cref{ex:sigma-assoc}} \\
  & \eqv{}{} \sm{a:A}{u:P(a)}{p:a=x}\trans{p}{f(a,u)}=v
  \tag{by \cref{thm:path-sigma}} \\
  & \eqv{}{} \sm{a:A}{p:a=x}{u:P(a)}\trans{p}{f(a,u)}=v
  \notag \\
  & \eqv{}{} \sm{u:P(x)}f(x,u)=v
  \tag{$*$}\label{eq:uses-sum-over-paths} \\
  & \jdeq \hfib{f(x)}{v}. \notag
\end{align}
  等价~\eqref{eq:uses-sum-over-paths} 由 \cref{thm:omit-contr,thm:contr-paths,ex:sigma-assoc} 得出。
\end{proof}

我们称一个逐纤维变换 $f:\prd{x:A} P(x)\to Q(x)$ 为\define{逐纤维等价}%
\indexdef{fiberwise!equivalence}%
\indexdef{equivalence!fiberwise}
如果每个 $f(x):P(x) \to Q(x)$ 都是一个等价。

\begin{thm}\label{thm:total-fiber-equiv}
假设 $f$ 是类型 $A$ 上两个族 $P$ 和 $Q$ 之间的一个逐纤维变换。
那么 $f$ 是逐纤维等价当且仅当 $\total{f}$ 是一个等价。
\end{thm}

\begin{proof}
令 $f$、$P$、$Q$ 和 $A$ 如定理所述。
由 \cref{fibwise-fiber-total-fiber-equiv} 可知，对所有的 $x:A$ 和 $v:Q(x)$，$\hfib{\total{f}}{\pairr{x,v}}$ 可缩当且仅当 $\hfib{f(x)}{v}$ 可缩。
因此，对所有 $w:\sm{x:A}Q(x)$ 而言 $\hfib{\total{f}}{w}$ 可缩，当且仅当对所有 $x:A$ 和 $v:Q(x)$ 而言 $\hfib{f(x)}{v}$ 可缩。
\end{proof}

\index{fiberwise!transformation|)}%

\section{对象分类子}
\label{sec:object-classification}

在类型论中，\emph{类型族}是一个基本概念，即一个函数 $B:A\to\type$。
我们已经看到，这样的族在某种程度上类似于同伦论中的\emph{纤维化}，其中的纤维化即为投影 $\proj1:\sm{a:A} B(a) \to A$。
同伦论中的一个基本事实是：每个映射都等价于一个纤维化。
有了泛等性，我们也可以在类型论中证明同样的结论。

\begin{lem}\label{thm:fiber-of-a-fibration}
  对于任何类型族 $B:A\to\type$，$\proj1:\sm{x:A} B(x) \to A$ 在 $a:A$ 上的纤维等价于 $B(a)$：
  \[ \eqv{\hfib{\proj1}{a}}{B(a)} \]
\end{lem}

\begin{proof}
  我们有
  \begin{align*}
    \hfib{\proj1}{a} &\defeq \sm{u:\sm{x:A} B(x)} \proj1(u)=a\\
    &\eqvsym \sm{x:A}{b:B(x)} (x=a)\\
    &\eqvsym \sm{x:A}{p:x=a} B(x)\\
    &\eqvsym B(a)
  \end{align*}
  这里利用了相等类型的左泛性质。
\end{proof}

\begin{lem}\label{thm:total-space-of-the-fibers}
  对于任意函数 $f:A\to B$，我们有 $\eqv{A}{\sm{b:B}\hfib{f}{b}}$。
\end{lem}

\begin{proof}
  我们有
  \begin{align*}
    \sm{b:B}\hfib{f}{b} &\defeq \sm{b:B}{a:A} (f(a)=b)\\
    &\eqvsym \sm{a:A}{b:B} (f(a)=b)\\
    &\eqvsym A
  \end{align*}
  这里利用了 $\sm{b:B} (f(a)=b)$ 可缩这一事实。
\end{proof}

\begin{thm}\label{thm:nobject-classifier-appetizer}
对任意类型 $B$，存在等价
\begin{equation*}
\chi:\Parens{\sm{A:\type} (A\to B)}\eqvsym (B\to\type).
\end{equation*}
\end{thm}

\begin{proof}
我们需要构造拟逆
\begin{align*}
\chi & : \Parens{\sm{A:\type} (A\to B)}\to B\to\type\\
\psi & : (B\to\type)\to\Parens{\sm{A:\type} (A\to B)}.
\end{align*}
定义 $\chi$ 为 $\chi((A,f),b)\defeq\hfiber{f}b$，定义 $\psi$ 为 $\psi(P)\defeq\Pairr{(\sm{b:B} P(b)),\proj1}$。
现在需要验证 $\chi\circ\psi\htpy\idfunc{}$ 以及 $\psi\circ\chi \htpy\idfunc{}$。
\begin{enumerate}
\item 设 $P:B\to\type$。
  由 \cref{thm:fiber-of-a-fibration}，
$\hfiber{\proj1}{b}\eqvsym P(b)$ 对任意 $b:B$ 成立，因此立即得到 $P\htpy\chi(\psi(P))$。
\item 设 $f:A\to B$ 为一个函数。我们需要找到一条道路
\begin{equation*}
\Pairr{\tsm{b:B} \hfiber{f}b,\,\proj1}=\pairr{A,f}.
\end{equation*}
首先注意，由 \cref{thm:total-space-of-the-fibers}，我们有
$e:\sm{b:B} \hfiber{f}b\eqvsym A$ 其中 $e(b,a,p)\defeq a$ 且 $e^{-1}(a)
\defeq(f(a),a,\refl{f(a)})$。
根据 \cref{thm:path-sigma}，只需证明 $\trans{(\ua(e))}{\proj1} = f$。
但由泛等公理的计算规则以及 \eqref{eq:transport-arrow}，我们有 $\trans{(\ua(e))}{\proj1} = \proj1\circ e^{-1}$，而 $e^{-1}$ 的定义立即给出 $\proj1 \circ e^{-1} \jdeq f$。\qedhere
\end{enumerate}
\end{proof}

\noindent
\indexdef{object!classifier}%
\indexdef{classifier!object}%
\index{.infinity1-topos@$(\infty,1)$-topos}%
特别地，这意味着在高阶意象理论的意义下我们有一个\emph{对象分类子}。
回顾 \cref{def:pointedtype}，$\pointed\type$ 表示带点类型 $\sm{A:\type} A$ 构成的类型。

\begin{thm}\label{thm:object-classifier}
设 $f:A\to B$ 为一个函数。则图表
\begin{equation*}
  \vcenter{\xymatrix{
      A\ar[r]^-{\vartheta_f} \ar[d]_{f} &
      \pointed{\type}\ar[d]^{\proj1}\\
      B\ar[r]_{\chi_f} &
      \type
      }}
\end{equation*}
是一个\emph{拉回}\index{pullback}方块（参见 \cref{ex:pullback}）。
其中函数 $\vartheta_f$ 定义为
\begin{equation*}
 \lam{a} \pairr{\hfiber{f}{f(a)},\pairr{a,\refl{f(a)}}}.
\end{equation*}
\end{thm}

\begin{proof}
注意到我们有等价
\begin{align*}
A & \eqvsym \sm{b:B} \hfiber{f}b\\
& \eqvsym \sm{b:B}{X:\type}{p:\hfiber{f}b= X} X\\
& \eqvsym \sm{b:B}{X:\type}{x:X} \hfiber{f}b= X\\
& \eqvsym \sm{b:B}{Y:\pointed{\type}} \hfiber{f}b = \proj1 Y\\
& \jdeq B\times_{\type}\pointed{\type}
\end{align*}
由此得到一个复合等价 $e:A\eqvsym B\times_\type\pointed{\type}$。
我们可以逐步展示这个复合等价的作用如下：
\begin{align*}
a & \mapsto \pairr{f(a),\; \pairr{a,\refl{f(a)}}}\\
& \mapsto \pairr{f(a), \; \hfiber{f}{f(a)}, \; \refl{\hfiber{f}{f(a)}}, \; \pairr{a,\refl{f(a)}}}\\
& \mapsto \pairr{f(a), \; \hfiber{f}{f(a)}, \; \pairr{a,\refl{f(a)}}, \; \refl{\hfiber{f}{f(a)}}}\\
& \mapsto \pairr{f(a), \; \pairr{\hfiber{f}{f(a)}, \; \pairr{a,\refl{f(a)}}}, \; \refl{\hfiber{f}{f(a)}}}.
\end{align*}
因此，我们得到同伦 $f\htpy\proj1\circ e$ 和 $\vartheta_f\htpy \proj2\circ e$。
\end{proof}

\section{泛等性蕴含函数外延性}
\label{sec:univalence-implies-funext}

\index{function extensionality!proof from univalence}%
在本章的最后一节，我们将证明泛等公理蕴含函数外延性。因此，在本节中我们\emph{不使用}函数外延性公理。
证明分为两步。首先，在 \cref{uatowfe} 中我们将证明泛等公理蕴含一种弱形式的函数外延性，即下面 \cref{weakfunext} 所定义的弱函数外延性原理。
而弱函数外延性原理又能推出通常的函数外延性，且这一步推导不需要泛等公理（\cref{wfetofe}）。

\index{univalence axiom}%
设 $\type$ 为一个宇宙；我们会明确指出在何处假设它是泛等的。

\begin{defn}\label{weakfunext}
\define{弱函数外延性原理}
\indexdef{function extensionality!weak}%
断言对于任意类型 $A$ 上的任意类型族 $P:A\to\type$，存在一个函数
\begin{equation*}
\Parens{\prd{x:A}\iscontr(P(x))} \to\iscontr\Parens{\prd{x:A}P(x)}
\end{equation*}
\end{defn}

下面的引理在已有函数外延性时容易证明；但这里的要点是，它也可以单独从泛等公理推出，而不需额外假设函数外延性。

\begin{lem} \label{UA-eqv-hom-eqv}
假设 $\type$ 是泛等的，则对任意 $A,B,X:\type$ 和任意 $e:\eqv{A}{B}$，存在一个等价
\begin{equation*}
\eqv{(X\to A)}{(X\to B)}
\end{equation*}
其底层映射就是对 $e$ 的底层函数做后复合。
\end{lem}

\begin{proof}
  % Immediate by induction on $\eqv{}{}$ (see \cref{thm:equiv-induction}).
  如同在 \cref{lem:qinv-autohtpy} 的证明中，我们可以假设 $e = \idtoeqv(p)$，其中 $p:A=B$。
  然后通过道路归纳，可设 $p$ 为 $\refl{A}$，于是 $e = \idfunc[A]$。
  但此时与 $e$ 的后复合就是恒等映射，因此是一个等价。
\end{proof}

\begin{cor}\label{contrfamtotalpostcompequiv}
设 $P:A\to\type$ 为一族可缩类型，即 \narrowequation{\prd{x:A}\iscontr(P(x)).}
则投影 $\proj{1}:(\sm{x:A}P(x))\to A$ 是一个等价。假设 $\type$ 是泛等的，与 $\proj{1}$ 的后复合立即给出一个等价
\begin{equation*}
\alpha : \eqv{\Parens{A\to\sm{x:A}P(x)}}{(A\to A)}.
\end{equation*}
\end{cor}

\begin{proof}
  根据 \cref{thm:fiber-of-a-fibration}，对 $\proj{1}:(\sm{x:A}P(x))\to A$ 和 $x:A$ 我们有等价
  \begin{equation*}
    \eqv{\hfiber{\proj{1}}{x}}{P(x)}.
  \end{equation*}
  因此当每个 $P(x)$ 可缩时，$\proj{1}$ 是一个等价。断言现在是 \cref{UA-eqv-hom-eqv} 的推论。
\end{proof}

特别地，上述等价在 $\idfunc[A]$ 处的同伦纤维是可缩的。因此，要证明泛等公理蕴含弱函数外延性，我们只需证明依值函数类型 $\prd{x:A}P(x)$ 是 $\hfiber{\alpha}{\idfunc[A]}$ 的收缩核。

\begin{thm}\label{uatowfe}
在泛等宇宙 $\type$ 中，设 $P:A\to\type$ 是一族可缩类型，并令 $\alpha$ 为 \cref{contrfamtotalpostcompequiv} 中的函数。
则 $\prd{x:A}P(x)$ 是 $\hfiber{\alpha}{\idfunc[A]}$ 的收缩核。从而 $\prd{x:A}P(x)$ 是可缩的。换言之，泛等公理蕴含弱函数外延性原理。
\end{thm}

\begin{proof}
定义函数
\begin{align*}
  \varphi &: (\tprd{x:A}P(x))\to\hfiber{\alpha}{\idfunc[A]},\\
  \varphi(f) &\defeq (\lam{x} (x,f(x)),\refl{\idfunc[A]}),
\intertext{and}
  \psi &: \hfiber{\alpha}{\idfunc[A]}\to \tprd{x:A}P(x), \\
  \psi(g,p) &\defeq \lam{x} \trans {\happly (p,x)}{\proj{2} (g(x))}.
\end{align*}
则由依值函数类型的唯一性原理，$\psi(\varphi(f))=\lam{x} f(x)$ 就是 $f$。
\end{proof}

我们现在来证明，弱函数外延性蕴含通常的函数外延性。
回想 \eqref{eq:happly} 中的函数 $\happly (f,g) : (f = g)\to(f\htpy g)$，它将函数的相等转换为同伦。
下面的证明不使用泛等公理。

\begin{thm}\label{wfetofe}
  \index{function extensionality}%
弱函数外延性蕴含函数外延性 \cref{axiom:funext}。
\end{thm}

\begin{proof}
我们想要证明
\begin{equation*}
\prd{A:\type}{P:A\to\type}{f,g:\prd{x:A}P(x)}\isequiv(\happly (f,g)).
\end{equation*}
根据 \cref{thm:total-fiber-equiv}，一个逐纤维映射诱导全空间上的等价，当且仅当它逐纤维为等价。因此，我们只需证明由 $\lam{g:\prd{x:A}P(x)} \happly (f,g)$ 诱导的、类型为
\begin{equation*}
\Parens{\sm{g:\prd{x:A}P(x)}(f= g)} \to \sm{g:\prd{x:A}P(x)}(f\htpy g)
\end{equation*}
的函数是等价。
由于 \cref{thm:contr-paths}，左边的类型是可缩的，故只需证明右边的类型：
\begin{equation}\label{eq:uatofesp}
\sm{g:\prd{x:A}P(x)}\prd{x:A}f(x)= g(x)
\end{equation}
是可缩的。
根据 \cref{thm:ttac}，该类型等价于
\begin{equation}\label{eq:uatofeps}
\prd{x:A}\sm{u:P(x)}f(x)= u.
\end{equation}
\cref{thm:ttac} 的证明虽然用到了函数外延性，但仅涉及其中的一个复合。
因此，在不假设函数外延性的前提下，我们仍可断言：~\eqref{eq:uatofesp} 是~\eqref{eq:uatofeps} 的一个收缩核\index{retract!of a type}。
而~\eqref{eq:uatofeps} 是可缩类型的积，根据弱函数外延性原理，它是可缩的；因此~\eqref{eq:uatofesp} 也是可缩的。
\end{proof}

\sectionNotes

在代数拓扑中，一个众所周知的事实是：配备了拟逆的连续映射空间，其同伦类型并非"同伦等价空间"应有的同伦类型。
在代数拓扑中，"同伦等价空间" $(\eqv ab)$ 通常被简单地定义为函数空间 $(A\to B)$ 的一个子空间，由那些作为同伦等价的函数组成。
在类型论中，与之最接近的是 $\sm{f:A\to B} \brck{\qinv(f)}$；参见 \cref{ex:brck-qinv}。

在同伦类型论中，等价的第一个定义是我们称为 $\iscontr(f)$ 的那个，它由 Voevodsky 提出。
其他定义随后被多人相继发现。
关于伴随等价\index{adjoint!equivalence}的基本定理，例如 \cref{lem:coh-equiv,thm:equiv-iso-adj}，是改编自高阶范畴论和同伦论中的标准事实。
将双可逆性用作等价的定义，是由 André Joyal 提出的。

\cref{sec:mono-surj,sec:equiv-closures} 中讨论的等价性质在同伦论中是众所周知的。
其中大多数在类型论中的首次证明由 Voevodsky 完成。

每个函数都等价于一个纤维化，这是同伦论中的一个标准事实。
$(\infty,1)$-范畴
\index{.infinity1-category@$(\infty,1)$-category}%
理论中的对象分类子
\index{object!classifier}%
\index{classifier!object}%
概念（即 \cref{thm:nobject-classifier-appetizer} 的范畴论类比）是由 Rezk（雷兹克）提出的（参见~\cite{Rezk05,lurie:higher-topoi}）。

最后，泛等公理蕴含函数外延性（\cref{sec:univalence-implies-funext}）这一事实归功于 Voevodsky。
我们给出的证明是他的证明的简化版本。
\cref{ex:funext-from-nondep} 同样归功于 Voevodsky。

\sectionExercises

\begin{ex}\label{ex:two-sided-adjoint-equivalences}
  考虑函数 $f:A\to B$ 的"双侧伴随等价\index{adjoint!equivalence}数据"所构成的类型，
  \begin{narrowmultline*}
    \sm{g:B\to A}{\eta: g \circ f \htpy \idfunc[A]}{\epsilon:f \circ g \htpy \idfunc[B]}
    \narrowbreak
    \Parens{\prd{x:A} \map{f}{\eta x} = \epsilon(fx)} \times
    \Parens{\prd{y:B} \map{g}{\epsilon y} = \eta(gy) }.
  \end{narrowmultline*}
  根据 \cref{lem:coh-equiv}，我们知道若 $f$ 是一个等价，则该类型是有实例的。
  试给出该类型的一个类似于 \cref{lem:qinv-autohtpy} 的刻画。

  你能否给出一个例子，说明该类型通常不是一个命题？
  （在 \cref{cha:hits} 之后这会更容易。）
\end{ex}

\begin{ex}\label{ex:symmetric-equiv}
  证明对于任意 $A,B:\UU$，下述类型等价于 $\eqv A B$。
  \begin{equation*}
    \sm{R:A\to B\to \type}
    \Parens{\prd{a:A} \iscontr\Parens{\sm{b:B} R(a,b)}} \times
    \Parens{\prd{b:B} \iscontr\Parens{\sm{a:A} R(a,b)}}.
  \end{equation*}
  你能否由此提取出一种满足 $\isequiv(f)$ 三项要求的类型定义？
\end{ex}

\begin{ex} \label{ex:qinv-autohtpy-no-univalence}
  试在不使用泛等公理的前提下，重新表述 \cref{lem:qinv-autohtpy} 的证明。
\end{ex}

\begin{ex}[不稳定的八面体公理]\label{ex:unstable-octahedron}
  \index{axiom!unstable octahedral}%
  \index{octahedral axiom, unstable}%
  假设 $f:A\to B$，$g:B\to C$，$b:B$。
  \begin{enumerate}
  \item 证明存在一个自然映射 $\hfib{g\circ f}{g(b)} \to \hfib{g}{g(b)}$，其在 $(b,\refl{g(b)})$ 上的纤维等价于 $\hfib f b$。
  \item 证明 $\eqv{\hfib{g\circ f}{c}}{\sm{w:\hfib{g}{c}} \hfib f {\proj1 w}}$。
  \end{enumerate}
\end{ex}

\begin{ex}\label{ex:2-out-of-6}
  \index{2-out-of-6 property}%
  证明等价满足\emph{2-out-of-6 性质（2-out-of-6 property）}：给定 $f:A\to B$，$g:B\to C$ 和 $h:C\to D$，如果 $g\circ f$ 和 $h\circ g$ 是等价，那么 $f$，$g$，$h$ 以及 $h\circ g\circ f$ 也都是等价。
  利用这一点给出 \cref{thm:paths-respects-equiv} 的一个更高层次的证明。
\end{ex}

\begin{ex}\label{ex:qinv-univalence}
  对于 $A,B:\UU$，以显然的方式通过道路归纳定义
  \[ \mathsf{idtoqinv}_{A,B} :(A=B) \to \sm{f:A\to B}\qinv(f) \]
  。
  令 \textbf{\textsf{qinv}-泛等性}表示泛等公理的一种修改形式，它断言对于所有 $A,B:\UU$，函数 $\mathsf{idtoqinv}_{A,B}$ 都有一个拟逆。
  \begin{enumerate}
  \item 证明在 \cref{sec:univalence-implies-funext} 中函数外延性的证明里，可以用 \qinv-泛等性代替泛等公理。
  \item 证明在 \cref{thm:qinv-notprop} 的证明里，可以用 \qinv-泛等性代替泛等公理。
  \item 证明 \qinv-泛等性是不一致的（即可以构造出 $\emptyt$ 的一个实例）。
    由此可见，在陈述泛等公理时，使用 $\isequiv$ 的"好"版本是至关重要的。
  \end{enumerate}
\end{ex}

\begin{ex}\label{ex:embedding-cancellable}
  证明一个函数 $f:A\to B$ 是嵌入，当且仅当以下两个条件成立：
  \begin{enumerate}
  \item $f$ 是\emph{左可消的}，即对任意 $x,y:A$，若 $f(x)=f(y)$ 则 $x=y$。\label{item:ex:ec1}
  \item 对任意 $x:A$，映射 $\apfunc f: \Omega(A,x) \to \Omega(B,f(x))$ 是一个等价。\label{item:ex:ec2}
  \end{enumerate}
  （特别地，若 $A$ 是一个集合，则 $f$ 是嵌入当且仅当它是左可消的，并且对所有的 $x:A$，$\Omega(B,f(x))$ 都是可缩的。）
  请举例说明条件~\ref{item:ex:ec1} 与条件~\ref{item:ex:ec2} 互不蕴含。
\end{ex}

\begin{ex}\label{ex:cancellable-from-bool}
  证明左可消函数 $\bool\to B$ 的类型（见 \cref{ex:embedding-cancellable}）等价于 $\sm{x,y:B}(x\neq y)$。
  类似地，请给出嵌入 $\bool\to B$ 的类型的一个显式刻画。
\end{ex}

\begin{ex}\label{ex:funext-from-nondep}
  \textbf{朴素非依值函数外延性公理}说的是：对于 $A,B:\type$ 和 $f,g:A\to B$，存在一个函数 $(\prd{x:A} f(x)=g(x)) \to (f=g)$。
  \indexdef{function extensionality!non-dependent}%
  修改 \cref{sec:univalence-implies-funext} 中的论证，以此证明该公理蕴含完整的函数外延性公理（\cref{axiom:funext}）。
\end{ex}

% Local Variables:
% TeX-master: "hott-online"
% End: